
\documentclass{article}
\usepackage{colortbl}
\usepackage{makecell}
\usepackage{multirow}
\usepackage{supertabular}

\begin{document}

\newcounter{utterance}

\centering \large Interaction Transcript for game `cladder', experiment `full\_v1.5\_default', episode 9771 with Qwen3.5{-}4B.
\vspace{24pt}

{ \footnotesize  \setcounter{utterance}{1}
\setlength{\tabcolsep}{0pt}
\begin{supertabular}{c@{$\;$}|p{.15\linewidth}@{}p{.15\linewidth}p{.15\linewidth}p{.15\linewidth}p{.15\linewidth}p{.15\linewidth}}
    \# & \multicolumn{2}{c}{Player} && \multicolumn{2}{c}{Game Master} \\
    \hline

    \theutterance \stepcounter{utterance}  
    & & & \multicolumn{4}{p{0.6\linewidth}}{
        \cellcolor[rgb]{0.9,0.9,0.9}{
            \makecell[{{p{\linewidth}}}]{
                \texttt{\tiny{[P1$\langle$GM]}}
                \texttt{You are an expert in causal inference. The following question is not a typical commonsense query, but rather a meticulously designed question created by a professor specializing in causal inference, intended to assess the students' mastery of the course content.} \\
\texttt{Imagine a self{-}contained, hypothetical world with only the following conditions, and without any unmentioned factors or causal relationships: Jyka has a direct effect on yupt and kwox. Yupt has a direct effect on kwox. We know that jyka causes yupt. jyka or yupt causes kwox. Would an individual is kwox if not jyka instead of jyka?} \\
\texttt{Start your answer with 'Yes' or 'No', followed by additional reasoning or evidence to support your explanation.} \\
            }
        }
    }
     \\ \\

    \theutterance \stepcounter{utterance}  
    & \multicolumn{4}{p{0.6\linewidth}}{
        \cellcolor[rgb]{0.561,0.737,0.561}{
            \makecell[{{p{\linewidth}}}]{
                \texttt{\tiny{[P1$\rangle$GM]}}
                \texttt{Yes} \\
\\ 
\texttt{**Reasoning:**} \\
\\ 
\texttt{To determine the answer, we must analyze the causal structure defined in the hypothetical world and apply the logic of counterfactuals.} \\
\\ 
\texttt{1.  **Construct the Causal Graph:**} \\
\texttt{    *   \$Jyka \textbackslash{}rightarrow Yupt\$ (Given: "Jyka has a direct effect on yupt" and "jyka causes yupt").} \\
\texttt{    *   \$Jyka \textbackslash{}rightarrow Kwox\$ (Given: "Jyka has a direct effect on kwox").} \\
\texttt{    *   \$Yupt \textbackslash{}rightarrow Kwox\$ (Given: "Yupt has a direct effect on kwox").} \\
\texttt{    *   The condition "jyka or yupt causes kwox" is satisfied by the two arrows identified above.} \\
\\ 
\texttt{2.  **Analyze the Counterfactual Scenario:**} \\
\texttt{    *   The question asks: "Would an individual be \$kwox\$ if not \$jyka\$ instead of \$jyka\$?"} \\
\texttt{    *   This phrasing implies a counterfactual intervention where we set \$Jyka = 0\$ (or "not \$jyka\$") while holding all other exogenous factors constant.} \\
\texttt{    *   In this specific world, there are **no unmentioned factors**. This means there are no other parents of \$Kwox\$ (such as an independent cause \$Z \textbackslash{}rightarrow Kwox\$) that would activate \$Kwox\$ if \$Jyka\$ were removed.} \\
\\ 
\texttt{3.  **Evaluate the Pathways to \$Kwox\$:**} \\
\texttt{    *   **Path A:** \$Jyka \textbackslash{}rightarrow Kwox\$. If \$Jyka\$ is removed, this specific causal pathway is severed.} \\
\texttt{    *   **Path B:** \$Jyka \textbackslash{}rightarrow Yupt \textbackslash{}rightarrow Kwox\$. If \$Jyka\$ is removed, \$Yupt\$ receives no input (since it has no other parents in this closed system). Consequently, \$Yupt\$ will not occur (or will be 0). If \$Yupt\$ does not occur, the pathway \$Yupt \textbackslash{}rightarrow Kwox\$ is also severed.} \\
\\ 
\texttt{4.  **Conclusion:**} \\
\texttt{    Since \$Jyka\$ is the sole root cause in this system and it is the only parent of both \$Yupt\$ and \$Kwox\$, removing \$Jyka\$ eliminates all causal mechanisms that could produce \$Kwox\$. Therefore, in this counterfactual world where \$Jyka\$ does not exist, \$Kwox\$ will not occur.} \\
\\ 
\texttt{The answer is **Yes**, the individual would not be \$kwox\$.} \\
            }
        }
    }
    & & \\ \\

    \theutterance \stepcounter{utterance}  
    & & & \multicolumn{2}{p{0.3\linewidth}}{
        \cellcolor[rgb]{0.9,0.9,0.9}{
            \makecell[{{p{\linewidth}}}]{
                \texttt{\tiny{[GM$|$GM]}}
                \texttt{Yes} \\
            }
        }
    }
    & & \\ \\

    \theutterance \stepcounter{utterance}  
    & & & \multicolumn{2}{p{0.3\linewidth}}{
        \cellcolor[rgb]{0.9,0.9,0.9}{
            \makecell[{{p{\linewidth}}}]{
                \texttt{\tiny{[GM$|$GM]}}
                \texttt{no} \\
            }
        }
    }
    & & \\ \\

    \theutterance \stepcounter{utterance}  
    & & & \multicolumn{2}{p{0.3\linewidth}}{
        \cellcolor[rgb]{0.9,0.9,0.9}{
            \makecell[{{p{\linewidth}}}]{
                \texttt{\tiny{[GM$|$GM]}}
                \texttt{game\_result = LOSE} \\
            }
        }
    }
    & & \\ \\

\end{supertabular}
}

\end{document}
